\documentclass[11pt]{article}

\usepackage[utf8]{inputenc}
\usepackage[T1]{fontenc}
\usepackage{amsmath, amssymb, bm}
\usepackage{geometry}
\geometry{margin=1.15in}

\title{Proximal NND Fitting with Known Heteroskedastic Errors and Unknown Homoskedastic Noise}
\author{}
\date{}

\begin{document}
\maketitle

\section{Model}

Let $\Omega \subseteq \{1,\ldots,m\}\times\{1,\ldots,n\}$ denote the observed
entries.  Suppose each observed entry has a known heteroskedastic standard error
$s_{ij}$, but that these known errors may miss an additional homoskedastic
variance component $g=\gamma^2$.  The likelihood is
\begin{equation}
    Y_{ij}\mid X_{ij}, g
        \sim \mathcal{N}\!\left(X_{ij},\, s_{ij}^2 + g\right),
        \qquad (i,j)\in\Omega .
\end{equation}
Missing entries are handled by omitting them from the likelihood.

We use a nuclear-norm distribution prior on $X$.  For fitting, it is most stable
to parameterize the grid directly in terms of the effective nuclear-norm rate
\[
    \eta = \lambda_{\mathrm{eff}},
\]
rather than a base rate $\lambda$ coupled to $g$.  Thus
\begin{equation}
    X\mid \eta
        \sim \mathrm{NND}(\eta),
        \qquad
    p(X\mid \eta)
        \propto
        \eta^{mn}
        \exp\!\left\{
            -\eta\|X\|_*
        \right\}.
\end{equation}
If a base homoskedastic parameterization is desired for reporting, it can be
recovered after fitting as
\begin{equation}
    \lambda = \eta\sqrt{g}.
\end{equation}
Under the effective-rate parameterization, the NND normalizer contributes
$-mn\log\eta$ to the negative log posterior but does not add an
$(mn/2)\log g$ term.  This substantially reduces the feedback between the
regularization path and the variance update.

For $g$ we use an optional inverse-gamma prior,
\begin{equation}
    p(g) \propto g^{-(a_0+1)}\exp(-b_0/g),
\end{equation}
but the fitting strategy below only requires a one-dimensional prior penalty in
$g$.

\section{Taylor Approximation for the Variance Update}

Let $\mu$ denote the current matrix estimate.  Define the row-wise effective
variances
\begin{equation}
    v_{ij}(g) = s_{ij}^2 + g,
    \qquad
    p_{ij}(g) =
    \begin{cases}
        1/v_{ij}(g), & (i,j)\in\Omega,\\
        0, & (i,j)\notin\Omega,
    \end{cases}
\end{equation}
and the nuclear-norm curvature proxy
\begin{equation}
    B(\mu)
        =
        \left(\mu^\top \mu + \epsilon I\right)^{-1/2}.
\end{equation}
For fixed $(\mu,g,\eta)$, the Taylor row precision is
\begin{equation}
    A_i(\mu,g,\eta)
        =
        \mathrm{diag}\{p_{ij}(g)\}_{j=1}^n
        + \eta B(\mu).
\end{equation}
The Taylor approximation therefore defines a full row covariance
\begin{equation}
    \Sigma_i(\mu,g,\eta)
        =
        A_i(\mu,g,\eta)^{-1}.
\end{equation}
This covariance is dense in general because $B(\mu)$ is dense.  For entrywise
scores and lightweight diagnostics we usually retain only the marginal
variances,
\begin{equation}
    \mathrm{diag}(\Sigma_i)
        =
        \mathrm{diag}\!\left(A_i(\mu,g,\eta)^{-1}\right).
\end{equation}
This diagonal extraction is a storage and scoring approximation, not an
assumption that the row covariance itself is diagonal.

Holding $\mu$ fixed, the collapsed Taylor objective for updating $g$ is,
up to constants independent of $g$,
\begin{align}
    \mathcal{J}_g(g\mid \mu,\eta)
        &=
        \frac{1}{2}
        \sum_{(i,j)\in\Omega}
        \left[
            \log\{s_{ij}^2+g\}
            +
            \frac{(y_{ij}-\mu_{ij})^2}{s_{ij}^2+g}
        \right]
        + \frac{1}{2}\sum_{i=1}^m \log |A_i(\mu,g,\eta)|
        \notag\\
        &\quad
        + \eta\|\mu\|_*
        + (a_0+1)\log g
        + \frac{b_0}{g}.
        \label{eq:hetero-homo-g-objective}
\end{align}
The first line is the heteroskedastic Gaussian likelihood plus the Taylor
log-determinant correction.  The second line is the fixed-rate NND prior and
the optional inverse-gamma prior on $g$.  The $\eta\|\mu\|_*$ and
$-mn\log\eta$ NND-normalizer terms are constant with respect to $g$, but they
are retained when scoring a grid of $\eta$ values.


\subsection{Stochastic Hutchinson--CG Update for $g$}

The exact update of $g$ is still expensive because each evaluation of
\eqref{eq:hetero-homo-g-objective} or its derivative requires the term
\begin{equation}
    \frac{1}{2}\sum_{i=1}^m \log |A_i(\mu,g,\eta)|.
\end{equation}
Naively this costs one Cholesky factorization of an $n\times n$ matrix for each
row.  Since $g$ is scalar, a practical alternative is to update $\ell=\log g$
using a stochastic gradient.  The likelihood and variance-prior
parts have closed-form derivatives.  Only the log-determinant derivative is
estimated stochastically.

For a fixed row $i$,
\begin{equation}
    \frac{\partial}{\partial \ell}
    \log |A_i|
        =
        \mathrm{tr}\!\left[
            A_i^{-1}
            \frac{\partial A_i}{\partial \ell}
        \right],
        \qquad \ell=\log g .
\end{equation}
Using Hutchinson probes $z_{ir}$ with
$\mathbb{E}(z_{ir}z_{ir}^{\top})=I$, this trace can be estimated by
\begin{equation}
    \mathrm{tr}\!\left[
        A_i^{-1}
        \frac{\partial A_i}{\partial \ell}
    \right]
    \approx
    \frac{1}{R}\sum_{r=1}^R
    z_{ir}^{\top}
    A_i^{-1}
    \frac{\partial A_i}{\partial \ell}
    z_{ir}.
\end{equation}
For the effective-rate parameterization,
\begin{align}
    \frac{\partial p_{ij}}{\partial \ell}
        &=
        -\frac{g}{(s_{ij}^2+g)^2},
        \qquad (i,j)\in\Omega, \\
    \frac{\partial}{\partial \ell}
    \eta
        &=
        0.
\end{align}
Therefore
\begin{equation}
    \frac{\partial A_i}{\partial \ell}
        =
        \mathrm{diag}\!\left\{
            -\frac{g}{(s_{ij}^2+g)^2}
            \mathbf{1}\{(i,j)\in\Omega\}
        \right\}_{j=1}^n .
\end{equation}

Each Hutchinson term only requires solving a linear system
\begin{equation}
    A_i x_{ir}
        =
        \frac{\partial A_i}{\partial \ell}z_{ir}.
\end{equation}
These solves are performed with conjugate gradients using matrix-vector
products
\begin{equation}
    x \mapsto
    \mathrm{diag}\{p_{ij}(g)\}_{j=1}^n x
    +
    \eta B(\mu)x.
\end{equation}
The row/probe solves are independent and can be vectorized over both $i$ and
$r$.  In implementation we keep the Hutchinson probes fixed during the fit and
warm start each CG solve from the previous solution for the same row/probe.
Fixed probes reduce stochastic jitter in the one-dimensional trajectory for
$g$, while warm starts exploit the fact that successive $g$ values usually
change slowly.

The resulting stochastic update is
\begin{equation}
    \ell^{(t+1)}
        =
        \min\left\{
            \ell_{\max},
            \max\left[
                \ell_{\min},
                \ell^{(t)}
                -
                \rho_t \widehat{\nabla}_{\ell}
    \mathcal{J}_g(e^{\ell^{(t)}}\mid \mu,\eta)
            \right]
        \right\},
\end{equation}
so the log-variance stays in the configured interval
$[\ell_{\min},\ell_{\max}]$.  The gradient may also be clipped for stability.
This update is no longer a monotone backtracking step because the
log-determinant gradient is stochastic, but in benchmarks a small fixed
learning rate with a handful of probes closely tracked the exact Cholesky
update while substantially reducing wall-clock time.


\section{Efficient Proximal Fit for the Matrix}

The full Taylor objective would include the log-determinant term in the update
for $\mu$, which is costly.  For efficiency, we instead update $\mu$ using the
weighted nuclear-norm proximal problem
\begin{equation}
    \hat\mu(g,\eta)
        =
        \arg\min_X
        \left\{
            \frac{1}{2}
            \sum_{(i,j)\in\Omega}
            \frac{(y_{ij}-X_{ij})^2}{s_{ij}^2+g}
            +
            \eta\|X\|_*
        \right\}.
        \label{eq:hetero-homo-proximal-x}
\end{equation}
This subproblem is solved by FISTA or proximal gradient with singular-value
thresholding.  The Taylor approximation is then used only to update $g$ and to
score candidate $\eta$ values.  This is substantially cheaper than including
the Taylor log-determinant in every gradient step for $X$.

The alternating fixed-$\eta$ fitting strategy is:
\begin{enumerate}
    \item Initialize $g>0$.
    \item With $g$ fixed, solve \eqref{eq:hetero-homo-proximal-x} for $\mu$.
    \item With $\mu$ fixed, update $\ell=\log g$ by a one-dimensional
          backtracking step on \eqref{eq:hetero-homo-g-objective}.
    \item Repeat until $g$ and $\mu$ stabilize, or until a fixed outer-iteration
          budget is reached.
\end{enumerate}
The effective penalty returned by the fit is
\begin{equation}
    \hat\lambda_{\mathrm{eff}}
        =
        \eta,
\end{equation}
which is the quantity directly comparable to a known-variance proximal
nuclear-norm penalty.

\section{Selecting $\eta=\lambda_{\mathrm{eff}}$}

In practice, $\eta$ should be selected by a grid or other outer search.  For
each candidate $\eta_k$, fit $(\hat\mu_k,\hat g_k)$ using the fixed-$\eta$
alternating scheme above.  Candidate fits can then be scored by one of the
following criteria.

\paragraph{Entrywise cross-validation.}
Hold out a subset of observed entries and choose the $\eta_k$ with best
predictive error or predictive likelihood on the held-out entries.

\paragraph{Taylor collapsed objective.}
Choose the candidate with the smallest final approximate objective, including
the fixed-rate NND normalizer and the Taylor log-determinant correction.

\paragraph{Analytical LOO.}
Use the marginal variances $\mathrm{diag}(\hat\Sigma_i)$ extracted from the full
row covariances and evaluate the Gaussian leave-one-entry predictive density
with effective variances $s_{ij}^2+\hat g_k$.

The Taylor and LOO scores are model-internal criteria, while cross-validation
is an external predictive criterion.  The internal criteria are much cheaper
than nested CV when a fitted grid is already available, but their performance
should be checked against CV in regimes of interest.

\section{Warning on Joint Estimation of $g$ and a Base $\lambda$}

Joint conditional-mode updates for $(g,\lambda)$ are not recommended for this
approximation.  Under the scaled prior, the conditional mode for the base
parameter $\lambda$ has the schematic form
\begin{equation}
    \hat\lambda
        \approx
        \frac{mn + a_\lambda - 1}
             {\|\mu\|_*/\sqrt{g} + b_\lambda}.
\end{equation}
Thus increasing $g$ can increase the base $\lambda$, while a larger base
$\lambda$ can in turn affect the next $g$ update through both
$\lambda/\sqrt{g}$ and the prior normalizer.  Empirically, this feedback can
make $\lambda$ and $g$ chase each other and run to boundary values.

For this reason, the recommended workflow is to hold the effective penalty
$\eta$ fixed within each fit, estimate $g$ conditional on that $\eta$, and
select $\eta$ using a grid scored by CV, Taylor ELBO/objective, or LOO.  A base
rate $\lambda=\eta\sqrt{\hat g}$ can be reported afterward if desired.  Joint
base-lambda--variance updates should be treated only as a diagnostic or
experimental option.

\end{document}
